\documentclass[final,3p,times,twocolumn]{elsarticle}
%\documentclass[final,3p,times]{elsarticle}
% https://ctan.org/pkg/elsarticle     2018-07-05 
% \documentclass[preprint, review,3p,times,twocolumn]{elsarticle}
% \documentclass[preprint,12pt]{elsarticle}
% Use the option review to obtain double line spacing
% \documentclass[preprint,review,12pt,]{elsarticle}

%\usepackage{paralist, enumerate}
\usepackage[ruled,linesnumbered]{algorithm2e}
\usepackage{amssymb,amsmath,amsthm,bm}
\usepackage{geometry}
\geometry{left=1.8cm,right=1.8cm}

\usepackage{graphicx, subfigure}
%\usepackage{booktabs, multirow, array}
%\usepackage{algorithm, algorithmic}
\usepackage{url}
%\usepackage[switch, mathlines]{lineno}
%\linenumbers\setlength\linenumbersep{1mm}

%\usepackage{enumitem}
%\usepackage{lineno}
%\linenumbers
%\modulolinenumbers[2]

\renewcommand{\vec}[1]{\mathbf{#1}}

%\newtheorem{proposition}{Proposition}
\graphicspath{{figures-pdf/}}
%\newlength\singleimagewidth
%\setlength\singleimagewidth{\columnwidth}

%\usepackage{ulem}

\journal{Optics and Lasers in Engineering}
\usepackage{hyperref}

\begin{document}

\begin{frontmatter}

\title{Cryptanalyzing a bit-level image encryption algorithm based on chaotic maps}

\author[zsi,uestc]{Heping Wen\corref{corr}}
\ead{hepingwen2016@gmail.com}

%\author[uestc]{Chongfu Zhang}

%\author[xtu-cn]{Eric Yong Xie}

\cortext[corr]{Corresponding author.}

\address[zsi]{School of Electronic and Information Engineering, University of Electronic Science and Technology of China, Zhongshan Institute, \\ Zhongshan 528402, China}

\address[uestc]{School of Information and Communication Engineering, University of Electronic Science and Technology of China, Chengdu 611731, China}

\begin{abstract}

Recently, a bit-level chaotic image encryption algorithm (BCIEA) was proposed \cite{LuX2016OLE}, in which a diffusion-confusion structure with only one round is adopted. 
And its security performance mainly relies on the dynamic mechanisms introduced in the diffusion and confusion phases respectively.
However, it is found that the BCIEA is neither injective  nor surjective, which deviates from a basic principle for designing symmetric cryptosystems.
On the basis of a reasonable and slight correction with BCIEA, we proposed a chosen ciphertext attack method for breaking BCIEA.
Theoretical analysis and experimental results verify the effectiveness and efficiency of the proposed attack method.


\end{abstract}
\begin{keyword}
image encryption \sep cryptanalysis \sep bit-level \sep chaos.
\end{keyword}
\end{frontmatter}

\section{Introduction}
\label{sec:intro}

Nowadays, with the continuous rapid development of network communication technology, the security of storage and transmission of multimedia information, especially images and video, is becoming more and more interesting.
The image has the characteristics of strong correlation and high redundancy of adjacent pixels, thus the traditional cryptographic algorithms, such as AES, DES and IDEA, are not suitable.
To cope with this challenge, many different methodologies have been studied for exploiting more promising image encryption schemes. 
However, due to the lack of authoritative metrics, the security of these cryptographic algorithms is still worth exploring and researching.

In recent years, bit-level technique has been introduced into chaotic image encryption for improving its security and efficiency performance. In 2010, Ye \textit{et. al} \cite{YE2010PRL} coined an image scrambling encryption of pixel bit using chaotic map. In 2011, Zhu \textit{et. al} \cite{zhuzl2011infsci} created a chaotic image encryption scheme bases on bit-level technique. Yet, In 2014, Zhang \textit{et. al} reported that \cite{zhuzl2011infsci} is vulnerable to chosen plaintext attack. In 2016, Li \textit{et. al} pointed out that \cite{YE2010PRL} cannot resist neither known plaintext attack nor chosen plaintext attack. Moreover, In our latest study, we found that a bit-level image cipher proposed in \cite{Shafique2018epjp} can be broken by chosen plaintext attack \cite{wenh2019epjp}. Nevertheless, comparing with the design of encryption algorithms, the research on cryptanalysis is still relatively lacking as mentioned in \cite{wenhp2019entropy}. 

%Nowadays, many image encryption technology has been a hot-spot in the field of information security. In 2019, Ghadirli \textit{et. al} \cite{Ghadirli2019SP} reviewed the current status of color image encryption technology and classified its methodologies into ten categories, including chaos theory, bit plane decomposition, DNA coding, and so on. 
%Yet, comparing with the design of encryption algorithms, the research on cryptanalysis is still relatively lacking. In the same year, Li C. \textit{et. al} \cite{Ghadirli2019SP,cqli:meet:ISAP19} reported the latest image encryption schemes from the viewpoints of a cryptanalyst,and pointed out that it is worth investing more research.
%Our latest works also validates this view \cite{wenhp2019entropy,wenh2019epjp}. In the previous studies, we found that an image encryption scheme \cite{song2015entropy} based on DNA coding can neither resist chosen plaintext attack nor chosen cipher attack, and a bit-level image cipher \cite{Shafique2018epjp} also vulnerable to chosen plaintext attack.

In 2016, Lu \textit{et. al} \cite{LuX2016OLE} proposed a bit-level image encryption algorithm based on chaotic maps named BCIEA, and claimed that BCIEA is sufficiently safe owing to the statistically numerical analysis results and theoretically security analysis. 
More specially, compared with other similar algorithms, the security of BCIEA is mainly reflected in: First, the diffusion and confusion parts are both based on bit-level; Second, a cyclic shift operation associated with intermediate ciphertext information is embedded in the diffusion; Third, during the confusion, the intermediate ciphertext information is used to update the initial value of the chaotic system, thereby dynamically generating the coordinate index sequences for the permutation encryption.
Yet, after cryptanalysis, the dynamic encryption process in diffusion and confusion is indeed more difficult for known plaintext attacks and selective plaintext attacks. However, for a ciphertext attack with a stronger attack strength, the B algorithm is still insecure. Exactly it was found that: on one hand, a bug existing in the decryption of BCIEA; on the other hand, based on a small rationalization correction on BCIEA, a chosen ciphertext attack method is proposed for deciphering. 

The remainder of the paper is organized as follows. Sec.~\ref{sec2:BCIEA} reviews BCIEA briefly;
Sec.~\ref{sec3:decryption} points out a bug about the decryption of BCIEA and gives a small rationalization correction;
Then, Sec.~\ref{sec3:Cryptanalysis} presents a comprehensive cryptanalysis on BCIEA. 
Section~\ref{sec4:Experiments} verifies the effectiveness of the attack method. The last section concludes the paper.

\section{The cryptosystem under study}
\label{sec2:BCIEA}

This section firstly presents a bit-level technique and the adopted chaotic map, and then introduces BCIEA briefly.

\subsection{A bit-level technique}

As \cite{ZhouY2014SP}, binary bitplane decomposition (BBD) is one of three bitplane decomposition technologies. Also, BBD is adopted in BCIEA for decomposing a grayscale image $\bm{P}$ into 8 binary bitplanes, give as
\begin{equation}
\bm{P} = \sum\limits_{n = 1}^{8} {2^{n - 1}\bm{p}_n} = \bm{p}_1 + 2\bm{p}_2 + 2^2\bm{p}_3 + ... + 2^{7}\bm{p}_{8}
\label{eq:BBD}
\end{equation}
where $\bm{P} \in \mathbb{Z}_{256}$, and $\bm{p}_n \in \mathbb{Z}_2$. Conversely, binary bitplane composition (BBC) combines 8 bit planes into one grayscale image.

\subsection{The adopted chaotic map}

In BCIEA, a chaotic map is used twice to generate several keystreams for encryption. The piece-wise linear chaotic map (PWLCM) is mathematically modeled by
\begin{equation}
\begin{aligned}
x_i = F(x_{i-1},\eta)=	\begin{cases}
x_{i-1}/ \eta,              & 0<x_{i-1}<\eta \\
(x_{i-1}- \eta)/(0.5-\eta), & \eta \le x_{i-1}<0.5 \\
F(1-x_{i-1},\eta),          & 0.5 \le x_{i-1}<1 \\
\end{cases}
\end{aligned}
\label{eq:PWLCM}
\end{equation}
where $x$ is the state variable, $x_0 \in (0,1)$ is the initial value and $\eta \in (0,0.5)$ is the control parameter, respectively. 

\subsection{Description of BCIEA}

The block diagram of the encryption process of BCIEA is shown in Fig. \ref{fig:BCIEA_en}. As can be seen, the main components of BCIEA are introduced briefly as follows:

\begin{figure*}[!htb]
	\centering
	\includegraphics[width=12cm]{BCIEA_EN_machine}
	\caption{The block diagram for the encryption process of BCIEA.}
	\label{fig:BCIEA_en}	
\end{figure*}  

\begin{itemize}
	
	\item \emph{The chaotic map and the Secret Key}:
	
	In BCIEA, two different groups of  initial values and control parameters $(x_0,\eta_1,y_0,\eta_2)$ are served as the secret keys.
	
	\item \emph{Initialization}:
	With the initial parameter $\eta_1$ and the initial value $x_0$, get a sequence $\bm{X}_1$ by iterating Eq. (\ref{eq:PWLCM}) for $N_0+MN$ times and discard the former $N_0$ values, given as
			\begin{equation*}
			X_1(i) = \bmod(floor(x(i))\times 10^{14},256)\\
			\label{eq:PRNS}
			\end{equation*}
			where $i =1 \sim MN$, $floor$ is the function that takes an integer, and $X_1(i) \in \mathbb{Z}_{256}$. Then, two sequences $\bm{b}_1,\bm{b}_2$ of length $L$ are obtained from $\bm{X_1}$ by the method of binary bitplane decomposition. Here, $\bm{b}_1,\bm{b}_2 \in \mathbb{Z}_{2}$.
%	\begin{itemize}
%		\item 
%		\item 
%	\end{itemize}
	
	\item \emph{Pre-processing by BBD}:
	For the sake of generality, the encrypted object is an 8-bit grayscale image of size $\bm{P}$ of size $M \times N$ (height $\times$ width). Convert the plain image $\bm{P}$ of size $M \times N$ (height $\times$ width) into two 1D sequences $\bm{A}_1$ and $\bm{A}_2$ of length $4MN$. For convenience, let $L$ be $4MN$.
	
	\item \emph{Stage 1. diffusion}:
	
%		\begin{itemize}
%			\item  Step 1. Firstly calculate the sum of $\bm{A}_2$, represented as
%			\begin{equation}
%			sum_1 = \sum\limits_{i= 1}^{L} A_2(i)
%			\label{eq:sum1}
%			\end{equation}
%			where $L=4MN$. Then run a cyclic right shift operation on $\bm{A}_1$ with $sum_1$ bits to obtain $\bm{A}_{11}$.
%			
%			\item  Step 2. Get $\bm{B}_{1}$ by performing the bitwise XOR operations between $\bm{A}_{11}$, $\bm{A}_{2}$ and $\bm{b}_{1}$, given as
%			\begin{equation}
%			\left\{ \begin{array}{l}
%			B_1(1) = A_{11}(1) \oplus A_{11}(L) \oplus A_{2}(1) \oplus b_{1}(1)\\
%			B_1(i) = A_{11}(i) \oplus A_{11}(i-1) \oplus A_{2}(i) \oplus b_{1}(i)
%			\end{array} \right.
%			\label{eq:B1_diffusion}
%			\end{equation}
%			where $i=2 \sim L$.
%			
%			\item  Step 3. Like Step 1, firstly compute the sum values $sum_2$ of $\bm{B}_1$ as
%			\begin{equation}
%			sum_2 = \sum\limits_{i= 1}^{L} B_1(i)
%			\label{eq:sum2}
%			\end{equation}
%			Then , get $\bm{A}_{22}$ cyclic right shifting $\bm{A}_2$ with $sum_2$ bits.
%			
%			\item  Step 4. Similar to Step 2, obtain $\bm{B}_{2}$ from $\bm{A}_{22}$, $\bm{B}_{1}$ and $\bm{b}_{2}$, modeled by
%			\begin{equation}
%			\left\{ \begin{array}{l}
%			B_2(1) = A_{22}(1) \oplus A_{22}(L) \oplus B_{1}(1) \oplus b_{2}(1)\\
%			B_2(i) = A_{22}(i) \oplus A_{22}(i-1) \oplus B_{1}(i) \oplus b_{2}(i)
%			\end{array} \right.
%			\label{eq:B2_diffusion}
%			\end{equation}
%			where $i=2 \sim L$.			
%		\end{itemize}
	
	Step 1. Firstly calculate the sum of $\bm{A}_2$, represented as
	\begin{equation}
	sum_1 = \sum\limits_{i= 1}^{L} A_2(i)
	\label{eq:sum1}
	\end{equation}
	where $L=4MN$. Then run a cyclic right shift operation on $\bm{A}_1$ with $sum_1$ bits to obtain $\bm{A}_{11}$.
	
	Step 2. Get $\bm{B}_{1}$ by performing the bitwise XOR operations between $\bm{A}_{11}$, $\bm{A}_{2}$ and $\bm{b}_{1}$, given as
	\begin{equation}
	\left\{ \begin{array}{l}
	B_1(1) = A_{11}(1) \oplus A_{11}(L) \oplus A_{2}(1) \oplus b_{1}(1)\\
	B_1(i) = A_{11}(i) \oplus A_{11}(i-1) \oplus A_{2}(i) \oplus b_{1}(i)
	\end{array} \right.
	\label{eq:A11_to_B1}
	\end{equation}
	where $i=2 \sim L$.
	
	Step 3. Like Step 1, firstly compute the sum values $sum_2$ of $\bm{B}_1$ as
	\begin{equation}
	sum_2 = \sum\limits_{i= 1}^{L} B_1(i)
	\label{eq:sum2}
	\end{equation}
	Then , get $\bm{A}_{22}$ cyclic right shifting $\bm{A}_2$ with $sum_2$ bits.
	
	Step 4. Similar to Step 2, obtain $\bm{B}_{2}$ from $\bm{A}_{22}$, $\bm{B}_{1}$ and $\bm{b}_{2}$, modeled by
	\begin{equation}
	\left\{ \begin{array}{l}
	B_2(1) = A_{22}(1) \oplus A_{22}(L) \oplus B_{1}(1) \oplus b_{2}(1)\\
	B_2(i) = A_{22}(i) \oplus A_{22}(i-1) \oplus B_{1}(i) \oplus b_{2}(i)
	\end{array} \right.
	\label{eq:B2_diffusion}
	\end{equation}
	where $i=2 \sim L$.
%	
	\item \emph{Stage 2. Confusion}:
	
	Step 1. Firstly compute the sum values of $\bm{B}_1$ and $\bm{B}_1$ by
	\begin{equation}
	sum = \sum\limits_{i= 1}^{L} (B_1(i)+B_2(i))
	\label{eq:sum3}
	\end{equation}
	Further update the initial value of the chaotic map by
	\begin{equation}
	s_0 = \bmod(y_0+sum/L,1)\\
	\label{eq:update_y0}
	\end{equation}
	Then, with the initial parameter $\eta_2$ and the initial value $s_0$,  using Eq. (\ref{eq:PWLCM}) again, generate a chaotic sequences of length $2L$, and then process it into two integer sequences $\bm{Y}$ and $\bm{Z}$ both of length $L$. Here, $\bm{Y},\bm{Z} \in [1,L]$.
	
	Step 2. Use $\bm{Y}$ and $\bm{Z}$ as coordinate indexes, and then get $\bm{C}_1$ and $\bm{C}_2$ by confusing the positions of the elements in $\bm{B}_2$ and $\bm{B}_1$ respectively. Exactly, given as
	\begin{equation}
	\left\{ \begin{array}{l}
	C_1(i) = B_2(Y(i))\\
	C_2(i) = B_1(Z(i))\\
	\end{array} \right.
	\label{eq:confusion}
	\end{equation}
	where $i=1 \sim L$. 
	
	\item \emph{Post-processing by BBC}:
	Finally, get the cipher image $\bm{C}$ from $\bm{C}_1$ and $\bm{C}_2$ by the method of BPC.
\end{itemize}
As stated in \cite{LuX2016OLE}, the decryption process is the inverse of the encryption process, but the specific steps are not described. For more details, please refer to \cite{LuX2016OLE}. 

\section{On the decryption of BCIEA}
\label{sec3:decryption}

\subsection{Problems with the decryption}

As everyone knows, ensuring the decryption is available is a basic condition for good cryptosystems. In addition, decryption is an integral part of cryptanalysis. For this reason, we discuss the decryption part firstly.
The block diagram of the decryption process of BCIEA is shown in Fig. \ref{fig:BCIEA_de}. 

\begin{figure*}[!htb]
	\centering
	\includegraphics[width=12cm]{BCIEA_DE_machine}
	\caption{The block diagram for the decryption process of BCIEA.}
	\label{fig:BCIEA_de}	
\end{figure*}  

\begin{itemize}
	
	\item \emph{Initialization}:
	As with encryption, Generate the keystreams $\bm{b}_1,\bm{b}_2$ by the secret keys $(\eta_1,x_0)$.

	\item \emph{Pre-processing by BBD}: Get $\bm{C}_1$ and $\bm{C}_2$ from the cipher image $\bm{C}$ by BPD. 
	
	\item \emph{Stage 1. Anti-confusion}:
	
	Step 1. Firstly compute the sum values of $\bm{C}_1$ and $\bm{C}_2$ by
	\begin{equation}
	sum = \sum\limits_{i= 1}^{L} (C_1(i)+C_2(i))
	\label{eq:sum_C}
	\end{equation}
	Note that the sum of $\bm{B}_1$ and $\bm{B}_2$ is equal to the sum of $\bm{C}_1$ and $\bm{C}_2$ because the confusion only changes the positions of the elements.	Then, with the same secret keys $(\eta_2,y_0)$, update the initial value $y_0$ to $s_0$ and further obtain the two coordinate indexes sequences $\bm{Y}$ and $\bm{Z}$.
	
	Step 2. Restore $\bm{B}_1$ and $\bm{B}_2$ from $\bm{C}_2$ and $\bm{C}_1$ by the anti-confuse, given by
	\begin{equation}
	\left\{ \begin{array}{l}
	B_1(Z(i)) = C_2(i)\\
    B_2(Y(i)) = C_1(i)
	\end{array} \right.
	\label{eq:anti-confusion}
	\end{equation}
	where $i=1 \sim L$. 
	
	
	\item \emph{Stage 2. Anti-diffusion}:
	
	Step 1. Corresponding to Step 4 of the confusion, recover $\bm{A}_{22}$ from $\bm{B}_{2}$, $\bm{B}_{1}$ and $\bm{b}_{2}$, modeled by
	\begin{equation}
	\left\{ \begin{array}{l}
	A_{22}(1) \oplus A_{22}(L)  = B_2(1) \oplus B_{1}(1) \oplus b_{2}(1)\\
	A_{22}(i) \oplus A_{22}(i-1) = B_2(i) \oplus B_{1}(i) \oplus b_{2}(i)
	\end{array} \right.
	\label{eq:B2_to_A22}
	\end{equation}
	where $i=2 \sim L$.
	
	Step 2. Obtain $sum_2$ by Eq. (\ref{eq:sum2}), and then get $\bm{A}_{2}$ by doing cyclic left shift operation on $\bm{A}_{22}$ with $sum_2$ bits.
	
	Step 3. Recover $\bm{A}_{11}$ by performing the bitwise XOR operations between $\bm{B}_{1}$, $\bm{A}_{2}$ and $\bm{b}_{1}$, given as
	\begin{equation}
	\left\{ \begin{array}{l}
	A_{11}(1) \oplus A_{11}(L) = B_1(1) \oplus A_{2}(1) \oplus b_{1}(1)\\
	A_{11}(i) \oplus A_{11}(i-1) B_1(i) \oplus A_{2}(i) \oplus b_{1}(i)
	\end{array} \right.
	\label{eq:B1_to_A11}
	\end{equation}
	where $i=2 \sim L$.
	
	Step 4. Obtain $sum_1$ by Eq. (\ref{eq:sum1}), and then get $\bm{A}_{1}$ by doing cyclic left shift operation on $\bm{A}_{11}$ with $sum_1$ bits.
	
	\item \emph{Post-processing by BBC}:
	Finally, get the plain image $\bm{A}$ from $\bm{A}_1$ and $\bm{A}_2$ by the method of BPC.
\end{itemize}

Observing Eqs. (\ref{eq:B2_to_A22}) and (\ref{eq:B1_to_A11}), the decryption results are the bitwise XOR values of two elements, instead of the values of an element. In fact, this paradigm will lead to solutions that are not unique. To this end, the following proof is given from the perspective of algebraic analysis. 

Take Eq. (\ref{eq:B2_to_A22}) as an example, let $u(i)=A_{22}(1)$, $v(i)=B_2(i) \oplus B_{1}(i) \oplus b_{2}(i)$. Thus, Eq. (\ref{eq:B2_to_A22}) becomes
\begin{equation}
\left\{ \begin{array}{l}
u(1) \oplus u(L)  = v(1)\\
u(i) \oplus u(i-1) = v(i)
\end{array} \right.
\label{eq:two_solutions}
\end{equation}
where $i=2 \sim L$, $u(i),v(i) \in \mathbb{Z}_{2}$. Here, $\bm{u}$ is known, and $\bm{v}$ is to be determined.
One can conclude that $\bm{u}$ has two solutions regardless of the value of $\bm{v}$ due to the symmetry of two unknowns in this equation.

For example, given $L=3$, one has $\bm{u},\bm{v} \in \bm{U}$. The mapping process is shown in Fig. \ref{fig:U2V}(a). $\bm{v}$=``101'' is known. Then, both ``110'' and ``001'' are solutions of $\bm{u}$ in Eq. (\ref{eq:two_solutions}).

%\begin{figure}[!htb]
%	\centering
%	\includegraphics[width=4cm]{U2V}
%	\caption{An example of injection and surjection in BCIEA.}
%	\label{fig:U2V}	
%\end{figure}  

\begin{figure}[!htb]
	\centering
	\subfigure[Neither injection nor surjection]{\includegraphics[width=4cm]{U2V_v1}}
	\subfigure[Both injection and surjection]{\includegraphics[width=4cm]{U2V_v2}}	
	\caption{The illustration diagram of two examples about injection and surjection in BCIEA with $L=3$.}
	\label{fig:U2V}
\end{figure}

Obviously, this violates the condition that the encryption and decryption functions required by the cryptosystem are single shot.

\subsection{A slight correction on the decryption}
\label{ssec:correction}

To ensure correct decryption, and refer to other similar algorithm designs \cite{ozkaynak:review:ND18}, a reasonable correction is to  remove $A_{11}(L),A_{22}(L)$ to meet both injection and surjection. Thus, similarly given $L=3$, the mapping process is shown in Fig. \ref{fig:U2V}(b). 
Then, without changing the idea of the algorithm, the equations of encryption do not need to be modified, except for their first equation of Eq. (\ref{eq:B1_diffusion}) and Eq. (\ref{eq:B2_diffusion}). Exactly, revised as
\begin{equation}
B_1(1) = A_{11}(1) \oplus A_{2}(1) \oplus b_{1}(1)
\end{equation}
and 
\begin{equation}
B_2(1) = A_{22}(1) \oplus B_{1}(1) \oplus b_{2}(1)
\label{eq:corr_B2_diffusion}
\end{equation}
respectively. Accordingly in the decryption, the first equation of  Eqs. (\ref{eq:B2_to_A22}) and (\ref{eq:B1_to_A11}) become
\begin{equation}
A_{22}(1) = B_2(1) \oplus B_{1}(1) \oplus b_{2}(1)
\end{equation}
and
\begin{equation}
A_{11}(1) = B_1(1) \oplus A_{2}(1) \oplus b_{1}(1)
\label{eq:corr_B1_to_A11}
\end{equation}
This ensures the correctness of the decryption.

\section{Cryptanalysis}
\label{sec3:Cryptanalysis}

For the adversary, the goal is to restore the acquired cipher images to the original plain images. A common way is to get its real or equivalent key. In BCIEA, the secret keys are $(x_0,\eta_1,y_0,\eta_2,A_{11}(0),A_{22}(0))$ considering the correction for the decryption given in Sec. \ref{ssec:correction}. 
%More specially, $(x_0,\eta_1,A_{11}(0),A_{22}(0))$ are for diffusion and $(y_0,\eta_2)$ are for confusion respectively.
In fact, the difficulty of finding $(x_0,\eta_1,y_0,\eta_2)$ is quite large, even impossible, because the preproces operation as shown in Equation 1 is used. To do this, we turn our goal to get their equivalent keys, exactly $(\bm{b}_1,\bm{b}_2,A_{11}(0),A_{22}(0))$ for the diffusion and $(\bm{Y},\bm{Z})$ for the confusion respectively.



\subsection{Analysis on the diffusion part}
\label{ssec:break_diffusion}
%In the diffusion, $\bm{b}_1,\bm{b}_2$ generated by the chaotic map is used as the diffusion sequences for encryption. 
The generation process of the two diffusion sequences $\bm{b}_1,\bm{b}_2$ are independent of the plain images. Also, $A_{11}(0)$ and $A_{22}(0)$ are not related to plain images. Thus, $(\bm{b}_1,\bm{b}_2,A_{11}(0),A_{22}(0))$ can be considered as a permanent equivalent key of the diffusion under a given secret key. 
%In addition, referring to Sec. \ref{ssec:correction}, $A_{11}(0)$ and $A_{22}(0)$ are also part of the key. Therefore, 


Based on chosen ciphertext attack, one can choose the all-zero cipher image $\bm{C}^{(0)}$ and get its corresponding plain image $\bm{A}^{(0)}$. Then, one gets $\bm{B}_1^{(0)}=0$ and $\bm{B}_2^{(0)}=0$ because the confusion does not change the element value. Also, $\bm{A}_1^{(0)}$ and $\bm{A}_2^{(0)}$ are also known according to the basic assumption of the cipher is public. At this point, the original algorithm is actually a diffusion-only process.
Eq. (\ref{eq:B2_to_A22}) becomes

%\end{equation}
\begin{equation}
\begin{aligned}
A_{22}^{(0)}(1) & = A_{22}(0) \oplus B_2^{(0)}(1) \oplus B_{1}^{(0)}(1) \oplus b_{2}(1)\\
& = A_{22}(0) \oplus b_{2}(1)\\
\end{aligned}
\label{eq:C0_1}	
\end{equation}
\begin{equation}
\begin{aligned}
A_{22}^{(0)}(i) & = A_{22}^{(0)}(i-1) \oplus B_2^{(0)}(i) \oplus B_{1}^{(0)}(i) \oplus b_{2}(i)\\
& = A_{22}^{(0)}(i-1) \oplus b_{2}(i)\\
\end{aligned}
\label{eq:C0_i}	
\end{equation}
where $i=2 \sim L$. Here, one has $sum_2=0$ because $\bm{B}_1^{(0)}=0$. Thus, $\bm{A}_2^{(0)}=\bm{A}_{22}^{(0)}$. Moreover, $A_{22}(0)$ is fixed for a given key. Let $\hat{b}_{2}(1)=A_{22}(0) \oplus b_{2}(1)$, Eq. (\ref{eq:C0_1}) evolves to

\begin{equation*}
\hat{b}_{2}(1)=A_{22}(0) \oplus b_{2}(1)=A_{22}^{(0)}(1) \\
\end{equation*}
And Eq. (\ref{eq:C0_i}) evolves to
\begin{equation*}
b_{2}(i) = A_{22}^{(0)}(i) \oplus A_{22}^{(0)}(i-1) \\
\end{equation*}
where $i=2 \sim L$. Thus, one can determine the result of $A_{22}(0) \oplus b_{2}(1)$, namely $\hat{b}_{2}(1)$, and $b_{2}(i)$ for $i=2 \sim L$. Note that $A_{22}(0)$ and $b_{2}(1)$ always appear in pairs, so there is no need to ask for specific $A_{22}(0)$ and $b_{2}(1)$. Here, let $\bm{\hat{b}}_2$ be a sequence identical to $\bm{b}_2$ except that only the first element is different, exactly is $\hat{b}_{2}(1)=A_{22}(0) \oplus b_{2}(1)$.

Secondly, one gets $\bm{A}_{11}^{(0)}$ from $\bm{A}_{1}^{(0)}$ by doing cyclic right shift with $sum_1$ bits. Here, $sum_1$ associated with $\bm{A}_{2}^{(0)}$ and $\bm{A}_{1}^{(0)}$ are both known, thus $\bm{A}_{11}^{(0)}$ is determined. Further, similarly following Eq. (\ref{eq:corr_B1_to_A11}), 
%\begin{equation}
%\begin{aligned}
%A_{11}^{(0)}(1) & = A_{11}(0) \oplus B_1^{(0)}(1) \oplus A_{2}^{(0)}(1) \oplus b_{1}(1)\\
%& = A_{2}^{(0)}(1) \oplus A_{11}(0) \oplus b_{1}(1)\\
%\end{aligned}
%\label{eq:A11_1}	
%\end{equation}
Let $\hat{b}_{1}(1)=A_{11}(0) \oplus b_{1}(1)$, one has
\begin{equation*}
\hat{b}_{1}(1)=A_{11}^{(0)}(1) \oplus A_{2}^{(0)}(1) \\
\end{equation*}
When $i=2 \sim L$, one has
\begin{equation*}
b_{1}(i) = A_{11}^{(0)}(i) \oplus  A_{11}^{(0)}(i-1) \\
\end{equation*}
Similarly, let $\bm{\hat{b}}_1$ be a sequence identical to $\bm{b}_1$ except that only the first element is different, exactly is $\hat{b}_{1}(1)=A_{11}(0) \oplus b_{1}(1)$.

Consequently, another equivalent key $\bm{\hat{b}}_1,\bm{\hat{b}}_2$ are determined. They can be used to eliminate the diffusion part for any given images. This also laid the foundation for the entire cryptanalysis. As a result, these two sequences are different for different plaintext or ciphertext.

\subsection{Analysis on the confusion part}
\label{ssec:break_confusion}

Once the diffusion can be eliminated, BCIEA degenerates into a confusion-only/permutation-only algorithm. Referring to existing research, many permutation-only encryption algorithms are unable to defend against known plaintext and selective plaintext attacks.
However, the objects deciphered in these studies are static replacement processes that are unrelated to plaintext. These analytical methods may not be directly applicable.

Unlike the diffusion part, the confusion keystrems $\bm{Y},\bm{Z}$ are dynamic for the different plain or cipher images because of Eqs. (\ref{eq:sum_C}) and (\ref{eq:update_y0}). Thus, there is no global/permanent equivalent key for the confusion.
Nevertheless, by observing Eqs. (\ref{eq:sum_C}) and (\ref{eq:update_y0}), one can see that This dynamic attribute also has certain rules. 
More, according to the given ciphertext, the sum value can be obtained, thereby providing a basis for selecting the ciphertext. This is also an important reason for choosing a ciphertext attack.

The purpose of the attacker is to restore the original information of the target ciphertext. By Eq. (\ref{eq:sum_C}), the value of $sum$ for the target cipher image $\bm{C}$ can be calculated. Thus, The diffusion sequences $\bm{Y},\bm{Z}$ is identical under the premise of selecting ciphertexts having the same sum value. Then, the following analysis will select a ciphertext attack to select a special ciphertext with the same sum value as the given ciphertext.

For a simple example, given $L=8$ and $\bm{Y}=(8,4,5,7,2,1,6,3)$. One can choose the cipher images as
\begin{equation}
	\left\{ \begin{array}{l}
		\bm{C}_1^{(1)} = (0,1,0,1,0,1,0,1) \\
		\bm{C}_1^{(2)} = (0,0,1,1,0,0,1,1) \\
		\bm{C}_1^{(3)} = (0,0,0,0,1,1,1,1) \\
	\end{array} \right.
	\label{eq:L8_C1}
\end{equation}
Simultaneously, construct the other part $\bm{C}_2^{(1)},\bm{C}_2^{(2)},\bm{C}_2^{(3)}$ so that the sum of the chosen cipher images are all $sum$. Thus, one of the forms is
\begin{equation}
	\left\{ \begin{array}{l}
		\bm{C}_2^{(1)} = (1,1,1,1,1,0,0,0) \\
		\bm{C}_2^{(2)} = (1,1,1,1,1,0,0,0) \\
		\bm{C}_2^{(3)} = (1,1,1,1,1,0,0,0) \\
	\end{array} \right.
	\label{eq:L8_C2}
\end{equation}
Correspondingly, one gets $\bm{B}_2$ by performing anti-confusion on $\bm{C}_1$, given by
\begin{equation}
\left\{ \begin{array}{l}
\bm{B}_2^{(1)} = (0,1,0,1,0,1,0,1) \\
\bm{B}_2^{(2)} = (0,0,1,1,0,0,1,1) \\
\bm{B}_2^{(3)} = (0,0,0,0,1,1,1,1) \\
\end{array} \right.
\label{eq:L8_B2}
\end{equation}
In the confusion, the elements of $\bm{B}_1$ are only replaced by $\bm{C}_2$, and the elements of $\bm{B}_2$ are only replaced by $\bm{C}_1$. Thus, one can get the confusion sequence $\bm{Y}$ by comparing all the elements in Eq. (\ref{eq:L8_C1}) and Eq. (\ref{eq:L8_B2}).

Similarly, the cipher images are constructed by exchanging $\bm{C}_1$ in Eq. (\ref{eq:L8_C1}) and $\bm{C}_2$ in Eq. (\ref{eq:L8_C2}), thus the above method can be used to find the confusion sequence $\bm{Z}$.
%The reason is that  because   related to the intermediate ciphertext 
%are related to the intermediate ciphertext
%The generation process of the two diffusion sequences $\bm{b}_1,\bm{b}_2$ are independent of the plain images. Also, $A_{11}(0)$ and $A_{22}(0)$ are not related to plain images. Thus, $(\bm{b}_1,\bm{b}_2,A_{11}(0),A_{22}(0))$ can be considered as a permanent equivalent key of the diffusion under a given secret key. 

\subsection{The proposed attack method}
\label{ssec:break_BCIEA}

Based on the above, the diffusion and the confusion of BCIEA cannot resist a chosen ciphertext attack method. Thus, for a given cipher image $\bm{\hat{C}}$, the specific steps for breaking BCIEA by chosen ciphertext attack is given as follows:
\begin{itemize}
	\item Step 1. Based on the condition of chosen ciphertext attack, choose the all-zero cipher image and the $2 \left\lceil {\log_{2}(L)} \right\rceil$ special cipher images with the same $sum$ value as $\bm{\hat{C}}$ by the method constructed in Sec. \ref{ssec:break_confusion}, and then use the decryption machine of BCIEA to get their corresponding cipher images respectively.
	
	\item Step 2. As Sec. \ref{ssec:break_diffusion}, get the two diffusion keystreams $\bm{{b}}_1,\bm{{b}}_2$ with the all-zero chosen cipher image and the corresponding plain image. 

	\item Step 3. Obtain the two confusion sequences $\bm{Y}$ and $\bm{Z}$ by the method given in Sec. \ref{ssec:break_confusion}.
	
	\item Step 4. Recover the original image $\bm{\hat{A}}$ of a given cipher image $\bm{\hat{C}}$ with the equivalent keys $\bm{{b}}_1,\bm{{b}}_2$ and $\bm{Y}$ and $\bm{Z}$. 

\end{itemize}

Note that the chosen cipher images constructed in Step 1 may change for a given different cipher image. Nevertheless, the attack method is still available and effective. Therefore, BCIEA cannot resist chosen ciphertext attack. 

For the data complexity of the attack method, the number of cipher images required is $1+2 \left\lceil {\log_{2}(L)} \right\rceil$.
Moreover, in the terms of computational complexity, using the decryption machine is $\bm{O}(1+2 \left\lceil {\log_{2}(L)} \right\rceil)$, determining the two equivalent keys $(\bm{\hat{b}}_1,\bm{\hat{b}}_2)$ and $(\bm{Y},\bm{Z})$ are both $\bm{O}(2L)$, and restoring the original image is $\bm{O}(8L)$.
%
\section{Experimental simulations and discussion}
\label{sec4:Experiments}
%

\subsection{Experimental setup}
%\label{ssec:Experiments}
To verify the effectiveness of the proposed attack method, some experiments were carried out. The experimental environment is Matlab R2016a based on a personal computer. Same as those in \cite{LuX2016OLE}, experimental image is an 8-bit grayscale image ``Lenna'' of size $256 \times 256$. Moreover, some other images with different size and type selecting from MISC image database.
considering that the secret keys of the cryptosystem are not fixed and unknown for attackers, some different keys are selected for experimental simulations.

%\subsection{Breaking BCIEA for a matrix of size $ 2 \times 2$}


\subsection{Breaking BCIEA for the same image as  \cite{LuX2016OLE}}
As \cite{LuX2016OLE}, an 8-bit grayscale image ``Lenna'' of size $256 \times 256$
was adopted as the typical experimental object. Accordingly, the image is also served as an example for the breaking experiments. For the adversary, the task is to recover the original version of the given cipher image $\bm{\hat{C}}$, which is shown in Fig. \ref{fig:Lena_P_C}(b).

\begin{figure}[!htb]
	\centering
	\subfigure[$\bm{\hat{A}}$]{
\includegraphics[height=2cm]{Lena_c}}
	\subfigure[$\bm{\hat{C}}$]{\includegraphics[height=2cm]{Lena}}	
	\caption{The image ``Lenna'' of size $256 \times 256$}
	\label{fig:Lena_P_C}
\end{figure}

Firstly, by Step 1 of Sec. \ref{ssec:break_BCIEA}, the all-zero chosen cipher image $\bm{C}^{0}$ and the corresponding plain image $\bm{A}^{0}$ of size $256 \times 256$ are shown in Fig. \ref{fig:all_zero_P_C}(a)-(b) respectively. Hence, one gets the diffusion keystreams $\bm{\hat{b}}_1,\bm{\hat{b}}_2$.
Secondly, by calculating $\bm{\hat{C}}$, one has $sum=32765$. Thus, with the method of Step 2 in Sec. \ref{ssec:break_BCIEA}, two groups of the chosen cipher images with this same value are given in Fig. \ref{fig:18a_cipher_images}(a)-(r) and Fig. \ref{fig:18b_cipher_images}(a)-(r) respectively. 
For the sake of simplicity, the illustrations of A are not given because they look similar to Fig. \ref{fig:all_zero_P_C}(b), exactly are noise-like images. Thirdly, one can get the confusion sequences $\bm{Y}$ and $\bm{Z}$ by Step 3 of Sec. \ref{ssec:break_BCIEA}. Finally, by Step 4 of Sec. \ref{ssec:break_BCIEA}, the original image is restored as shown in Fig. \ref{fig:Lena_P_C}(a).

\begin{figure}[!htb]
	\centering
	\subfigure[$\bm{C}^{0}$]{\includegraphics[height=2cm]{all_255}}
	\subfigure[$\bm{A}^{0}$]{
\includegraphics[height=2cm]{cipher_all_255}}	
	\caption{The all-zero cipher image and the corresponding plain image.}
	\label{fig:all_zero_P_C}
\end{figure}


\begin{figure}[!htb]
	\centering
	\subfigure[$\bm{C}^{18}$]{
\includegraphics[height=1.3cm]{CC18a}}
	\subfigure[$\bm{C}^{17}$]{
\includegraphics[height=1.3cm]{CC17a}}
	\subfigure[$\bm{C}^{16}$]{
\includegraphics[height=1.3cm]{CC16a}}
	\subfigure[$\bm{C}^{15}$]{
\includegraphics[height=1.3cm]{CC15a}}
	\subfigure[$\bm{C}^{14}$]{
\includegraphics[height=1.3cm]{CC14a}}
	\subfigure[$\bm{C}^{13}$]{
\includegraphics[height=1.3cm]{CC13a}}\\
	\subfigure[$\bm{C}^{12}$]{
\includegraphics[height=1.3cm]{CC12a}}
	\subfigure[$\bm{C}^{11}$]{
\includegraphics[height=1.3cm]{CC11a}}
	\subfigure[$\bm{C}^{10}$]{
\includegraphics[height=1.3cm]{CC10a}}
	\subfigure[$\bm{C}^{9}$]{
\includegraphics[height=1.3cm]{CC9a}}
	\subfigure[$\bm{C}^{8}$]{
\includegraphics[height=1.3cm]{CC8a}}
	\subfigure[$\bm{C}^{7}$]{
\includegraphics[height=1.3cm]{CC7a}}\\
	\subfigure[$\bm{C}^{6}$]{
\includegraphics[height=1.3cm]{CC6a}}
	\subfigure[$\bm{C}^{5}$]{
\includegraphics[height=1.3cm]{CC5a}}
	\subfigure[$\bm{C}^{4}$]{
\includegraphics[height=1.3cm]{CC4a}}
	\subfigure[$\bm{C}^{3}$]{
\includegraphics[height=1.3cm]{CC3a}}
	\subfigure[$\bm{C}^{2}$]{
\includegraphics[height=1.3cm]{CC2a}}
	\subfigure[$\bm{C}^{1}$]{
\includegraphics[height=1.3cm]{CC1a}}		
	\caption{The 18 chosen cipher images for determining $\bm{Y}$.}
	\label{fig:18a_cipher_images}
\end{figure}

\begin{figure}[!htb]
	\centering
	\subfigure[$\bm{C}^{18}$]{
\includegraphics[height=1.3cm]{CC18b}}
	\subfigure[$\bm{C}^{17}$]{
\includegraphics[height=1.3cm]{CC17b}}
	\subfigure[$\bm{C}^{16}$]{
\includegraphics[height=1.3cm]{CC16b}}
	\subfigure[$\bm{C}^{15}$]{
\includegraphics[height=1.3cm]{CC15b}}
	\subfigure[$\bm{C}^{14}$]{
\includegraphics[height=1.3cm]{CC14b}}
	\subfigure[$\bm{C}^{13}$]{
\includegraphics[height=1.3cm]{CC13b}}\\
	\subfigure[$\bm{C}^{12}$]{
\includegraphics[height=1.3cm]{CC12b}}
	\subfigure[$\bm{C}^{11}$]{
\includegraphics[height=1.3cm]{CC11b}}
	\subfigure[$\bm{C}^{10}$]{
\includegraphics[height=1.3cm]{CC10b}}
	\subfigure[$\bm{C}^{9}$]{
\includegraphics[height=1.3cm]{CC9b}}
	\subfigure[$\bm{C}^{8}$]{
\includegraphics[height=1.3cm]{CC8b}}
	\subfigure[$\bm{C}^{7}$]{
\includegraphics[height=1.3cm]{CC7b}}\\
	\subfigure[$\bm{C}^{6}$]{
\includegraphics[height=1.3cm]{CC6b}}
	\subfigure[$\bm{C}^{5}$]{
\includegraphics[height=1.3cm]{CC5b}}
	\subfigure[$\bm{C}^{4}$]{
\includegraphics[height=1.3cm]{CC4b}}
	\subfigure[$\bm{C}^{3}$]{
\includegraphics[height=1.3cm]{CC3b}}
	\subfigure[$\bm{C}^{2}$]{
\includegraphics[height=1.3cm]{CC2b}}
	\subfigure[$\bm{C}^{1}$]{
\includegraphics[height=1.3cm]{CC1b}}		
	\caption{The 18 chosen cipher images for determining $\bm{Y}$.}
	\label{fig:18b_cipher_images}
\end{figure}


For the complexity of the attack method, the number of the chosen cipher images required is 37, and the total breaking time is about 3.27s. As can be seen, the attack method is both effective and efficient for breaking BCIEA.


\subsection{Breaking BCIEA for other images}

To further confirm the effectiveness of the attack method, some other images selected from the MISC database were also tested. Fig. \ref{fig:exp_results}(a)-(f) show the attacking results for the three images. Table \ref{tab:exp_results} gives the results for attacking BCIEA with some different images. It can be seen that, for these images with different types and sizes, the attack method is valid, and the required time cost is low. 


\begin{figure}[!htb]
	\centering
	\subfigure[]{
\includegraphics[height=2cm]{grayCamera_c}}
	\subfigure[]{\includegraphics[height=2cm]{grayCamera}} 
	\subfigure[]{
\includegraphics[height=2cm]{misc512A_c}}
	\subfigure[]{\includegraphics[height=2cm]{misc512A}}
	\caption{Experimental results for breaking BCIEA.}
	\label{fig:exp_results}	
\end{figure}

%

\begin{table}[!htb]
	\label{tab:exp_result}
	\caption{Attacking results for breaking BCIEA.}
	\centering
	\begin{tabular}{ccc}
		%		\toprule
		\hline
		\textbf{Images}  & \textbf{Encryption/Breaking times} & \textbf{Complexity} \\
		\hline 
		Lenna & 0.0647s / 0.2539s  & 37  \\ 
		Cameramen & U 0.0647s / 0.2539s  & 41  \\  
		Baboon& 0.0647s / 0.2539s  & 37  \\ 
%		Gray Lenna of size 256 $\times$ 256 & 0.0647s / 0.2539s  & 37  \\ 
%		Gray Cameramenof size 512 $\times$ 512 & U 0.0647s / 0.2539s  & 41  \\  
%		Color Baboon of size 256 $\times$ 256 & 0.0647s / 0.2539s  & 37  \\ 
		%		ucid00416 & RGB color 512 $\times$ 384 & UCID & 0.0647s & 0.2539s  & 2  \\ 
		%		ucid00622 & RGB color 384 $\times$ 512 & UCID & 0.0647s & 0.2539s  & 2  \\ 
		\hline
	\end{tabular}
	\label{tab:exp_results}
\end{table}



\section{Conclusions}


This paper performed a comprehensive cryptanalysis on a bit-level image encryption scheme based on chaotic maps so-called BCIEA. We found that BCIEA cannot resist a chosen ciphertext attack because of its inherent security defects. Both theoretical analysis and experimental verification were provided to support that BCIEA can be broken with low computational and data complexity. This study can provide a reference for improving the security of image encryption scheme based on bit-level technique.


\cite{Shafique2018epjp,MPEanalysis2015infsci,Diaconu2015infsci,Stinson2006,fengw2018ieeepj,wenh2019epjp,zhuzl2011infsci,Lcq:Optimal:SP11,Jolfaei2016TransIFS,Uhl2018Trans,YE2010PRL,fridrich1998symmetric,lic2016sigpro,ozkaynak:review:ND18,shannon1949communication,tengl2017MTA}

\section*{Acknowledgement}

%This research was supported by the Natural Science Foundation of China (No.~61772447, 61532020), the Joint Funds of the National Natural Science Foundation of China, China General Technology Research Institute (No.~U1736113) and the Fundamental Research Funds for the Central Universities (No.~531118010348).

\bibliographystyle{elsarticle-num}
\bibliography{OLE_BIEA,IEC-BPMC}
\end{document}